\scanpage{663}
\mspage{634}

On va considérer aussi l'hom (en fait effectivement \unsure{mis en jeu} par (38), même dans des cas pas du tout dégénérés tels que $\unsure{\mathfrak{G}|\mathcal{G}}^+=\mathrm{SL}(2,\widehat{\mathbb Z})$)
\[
(45)\qquad
\diagram{\begin{tikzcd}[ampersand replacement=\&,column sep=large]
G_{\rho,\sigma}^{\unsure{\natural|\sharp}}\arrow[r]\&\operatorname{Aut}(\unsure{\mathfrak{G}|\mathcal{G}}^+)\times\mu\times\mu\\
\underline u\arrow[r,mapsto]\&(\operatorname{int}(\underline u)|\unsure{\mathfrak{G}|\mathcal{G}}^+,\varepsilon_3(\underline u),\varepsilon_2(\underline u))
\end{tikzcd}}
\]
induit par (38), son image est formée des systèmes $(u,\varepsilon_3,\varepsilon_2)$ tels que
\[
(46)\qquad
\left\{\begin{array}{ll}
\text{a)}&u(\rho)\text{ est }\unsure{\mathfrak{G}|\mathcal{G}}^{\unsure{\natural|\sharp}}\text{-conjugué à }\varepsilon_3[\tau](\rho)\\
\text{b)}&u(\sigma)\quad\text{\rule{4.5cm}{0.3pt}}\quad\varepsilon_2[\tau](\sigma)\\
\text{c)}&u(L_0^{\unsure{\mathfrak{G}|\mathcal{G}}})\text{ est }\unsure{\mathfrak{G}|\mathcal{G}}^{\unsure{\natural|\sharp}}\text{-conjugué à }L_0\\
\text{d)}&u\text{ stabilise }\unsure{\mathfrak{G}|\mathcal{G}}^{\unsure{\natural|\sharp}}\quad\text{(automatique si }\varepsilon_3=\varepsilon_2,\ \text{\unsure{vu a) b)|vu c) b)}})
\end{array}\right.
\]
Son noyau est
\[
(47)\quad
\left\{f^{-1}\underline u\ \unsure{\in|\notin}\ G_{\rho,\sigma}^{\unsure{\natural|\sharp}}\ \middle|\
\begin{array}{l}\underline u=(u,\mu,1,1)\in Z_{\rho,\sigma}^{\circ}\\ f\in\unsure{\mathfrak{G}|\mathcal{G}}^{\unsure{\natural|\sharp}}\end{array}
\quad\text{\unsure{t.q.}}\quad u=\operatorname{int}(f)\right\}.
\]

Pour apprécier la structure du noyau, soit
\[
(48)\quad
\begin{aligned}
SN_0^{\deleted{\unsure{\mathfrak G}}}&=\operatorname{Norm}_{\unsure{\mathfrak{G}|\mathcal{G}}^+}(L_0),\qquad
SN_0^{\deleted{\unsure{\mathfrak G}}\unsure{\natural|\sharp}}=SN_0\cap\unsure{\mathfrak{G}|\mathcal{G}}^{\unsure{\natural|\sharp}}\\
&\hspace{3cm}=\left\{f\in\unsure{\mathfrak{G}|\mathcal{G}}^{\unsure{\natural|\sharp}}\ \middle|\ \exists\mu\in\widehat{\mathbb Z}^{*}\ \text{\unsure{t.q.}}\
\underbrace{(\operatorname{int}f,\mu,1,1)}_{\in Z_{\rho,\sigma}}\right\}
\end{aligned}
\]
et soit
\[
(49)\quad
\begin{gathered}
S\widetilde N_0\subset SN_0\times\widehat{\mathbb Z}^{*},\qquad
\unsure{S\widetilde N_0^{\unsure{\natural|\sharp}}|\widetilde N_0^{\unsure{\natural|\sharp}}}
=S\widetilde N_0\cap(N_0^{\unsure{\natural|\sharp}}\times\widehat{\mathbb Z}^{*})\\
\begin{array}{c}\Vert\\ \{(u,\mu)\mid u(\xi_0)=\xi_0^{\mu}\}\end{array}
\end{gathered}
\]
\margin{$\unsure{\widetilde N_0^{\unsure{\natural|\sharp}}}$}
--- donc on a
\[
(50)\quad
\diagram{\begin{tikzcd}[ampersand replacement=\&]
S\widetilde N_0^{\unsure{\natural|\sharp}}\arrow[r,"\sim"]\&SN_0^{\unsure{\natural|\sharp}}\\
(u,\mu)\arrow[r,mapsto]\&u
\end{tikzcd}}
\quad\text{si}\quad
\diagram{\begin{tikzcd}[ampersand replacement=\&]
\widehat{\mathbb Z}\arrow[r]\&L_0^{\unsure{\mathfrak{G}|\mathcal{G}}}\\
n\arrow[r,mapsto]\&\xi_0^n
\end{tikzcd}}
\quad\text{est injectif.}
\]
\margin{idem pour $S\widetilde N_0$.}
On a un hom canonique injectif
\[
(51)\quad
\diagram{\begin{tikzcd}[ampersand replacement=\&,column sep=large]
S\widetilde N_0^{\unsure{\natural|\sharp}}\arrow[r]\&S_0\widetilde G_{\rho,\sigma}^{\unsure{\natural|\sharp}}=Z_{\rho,\sigma}\cdot\unsure{\mathfrak{G}|\mathcal{G}}^+\\
(f,\mu)\arrow[r,mapsto]\&f^{-1}\cdot\underbrace{(\operatorname{int}(f),\mu,1,1)}_{\unsure{u|\underline u}}=\unsure{\underline u|u}\cdot f^{-1}
\end{tikzcd}}
\]
\margin{662}


\scanpage{664}

qui induit sur $L_0^{\unsure{\natural|\sharp}}$ l'hom canonique \inserted{(28)} (et qui prolonge (42)) --- il est encore vrai que l'image inverse \deleted{\unsure{de}} de \deleted{$\unsure{L_0^{\natural}}$} $i_{\unsure{L_0,G}}(L_0^{\unsure{\natural|\sharp}})$ est $L_0^{\unsure{\natural|\sharp}}\subset S\widetilde N_0^{\unsure{\natural|\sharp}}$, et on trouve ainsi une suite exacte canonique.
\[
(\unsure{52|51})\qquad
\diagram{\begin{tikzcd}[ampersand replacement=\&,column sep=small]
1\arrow[r,no head]\&L_0^{\unsure{\natural|\sharp}}\arrow[r]\&S\widetilde N_0^{\unsure{\natural|\sharp}}\arrow[r]\&G_{\rho,\sigma}^{\unsure{\natural|\sharp}}\arrow[r]\&\widetilde{\widetilde G}_{\rho,\sigma}\arrow[r,no head]\arrow[d,hook,no head]\&1\\
\&\&\&\&\operatorname{Aut}(\unsure{\mathfrak{G}|\mathcal{G}}^+)\times\mu\times\mu\&
\end{tikzcd}}
\]
où \deleted{\unsure{avec} $\unsure{\mathfrak{G}|\mathcal{G}}$} $\widetilde{\widetilde G}_{\rho,\sigma}$ est formé des $(u,\varepsilon_2,\varepsilon_3)$ satisfaisant (\deleted{\unsure{4}}46) --- i.e. on a
\[
(52\text{ bis})\qquad
S\widetilde N_0^{\unsure{\natural|\sharp}}/L_0^{\unsure{\natural|\sharp}}\simeq
\operatorname{Ker}\bigl(G_{\rho,\sigma}^{\unsure{\natural|\sharp}}\longrightarrow\operatorname{Aut}(\unsure{\mathfrak{G}|\mathcal{G}}^+)\times\mu\times\mu\unsure{;}\bigr).
\]
\margin{663}


\scanpage{665}
\mspage{635}

$4^\circ)$ $N_\rho,N_\sigma,Z_\rho,Z_\sigma,N_\rho^{\unsure{\natural|\sharp}},N_\sigma^{\unsure{\natural|\sharp}},Z_\rho^{\unsure{\natural|\sharp}},Z_\sigma^{\unsure{\natural|\sharp}}$.

On pose
\[
(53)\quad
\left\{\begin{array}{l}
N_\rho=\{\unsure{\mathbf a|a}\in G_{\rho,\sigma}\mid \unsure{\mathbf a|a}(\rho)=\rho^{\varepsilon_3},\ \text{où }\varepsilon_3=\varepsilon_3(a)\}\\
N_\sigma=\{b\in G_{\rho,\sigma}\mid b(\unsure{\rho|\sigma})=\unsure{\rho|\sigma}^{\varepsilon_2},\ \text{\unsure{où}}\ \varepsilon_2=\varepsilon_2(b)\}
\end{array}\right.
\]
\uline{NB} Il n'y a pas lieu finalement de \unsure{traîner|trainer} des $*$ en donnant une définition des $N_\rho,N_\sigma$ qui donnerait des groupes plus gros. (\unsure{Dans} p 580, \unsure{580'} je me suis \unsure{entrecorné|entortillé}\deleted{$\unsure{\natural}$} à ce sujet...)

On pose
\[
(54)\quad
\left\{\begin{aligned}
Z_\rho&=\operatorname{Cent}_{G_{\rho,\sigma}}(\rho)=\operatorname{Cent}_{G_{\rho,\sigma}}(L_\rho)
=\operatorname{Ker}(N_\rho\xrightarrow{\varepsilon_3}\mu).\\
Z_\sigma&=\operatorname{Cent}_{G_{\rho,\sigma}}(\sigma)=\operatorname{Cent}_{G_{\rho,\sigma}}(L_\sigma)
=\operatorname{Ker}(N_\sigma\xrightarrow{\varepsilon_2}\mu)
\end{aligned}\right.
\]
\[
(55)\quad
\left\{\begin{aligned}
SZ_\rho&=Z_\rho\cap\unsure{\mathfrak{G}|\mathcal{G}}^+=\operatorname{Ker}(Z_\rho\xrightarrow{\delta_Z}\gamma_{\rho,\sigma})\\
&=\operatorname{Cent}_{\unsure{\mathfrak{G}|\mathcal{G}}^+}(\rho)=\operatorname{Cent}_{\unsure{\mathfrak{G}|\mathcal{G}}^+}(L_\rho)\\
SZ_\sigma&=Z_\sigma\cap\unsure{\mathfrak{G}|\mathcal{G}}^+=\operatorname{Ker}(Z_\sigma\xrightarrow{\delta_{\unsure{Z|z}}}\gamma_{\rho,\sigma})\\
&=\operatorname{Cent}_{\unsure{\mathfrak{G}|\mathcal{G}}^+}(\sigma)=\operatorname{Cent}_{\unsure{\mathfrak{G}|\mathcal{G}}^+}(L_\sigma)
\end{aligned}\right.,
\]
\margin{on définit de façon identique $SN_\rho,SN_\sigma$, mais on a $SN_\rho=SZ_\rho$, $SN_\sigma=SZ_\sigma$.}
d'où
\[
(56)\quad
\diagram{\begin{tikzcd}[ampersand replacement=\&]
1\arrow[r]\&\unsure{Z_\rho}\arrow[r]\&N_\rho\arrow[r,"\varepsilon_{3\rho}"]\&\mu\arrow[r]\&\deleted{\unsure{\text{surjectif}}}\\
1\arrow[r,no head]\&Z_\sigma\arrow[r,no head]\&N_\sigma\arrow[r,"\varepsilon_{2\sigma}"]\&\mu\arrow[r,no head]\&1
\end{tikzcd}}
\]
\margin{Suites exactes \uline{semi-directes} par $\mu_{\tau'}$ --- la dernière aussi \unsure{semi-dir} par $\mu_\tau$.}
(Le fait qu'on peut mettre des $1$ au bout provient du fait que $\tau'\in N_\rho\cap N_\sigma$ satisfait
\[
(57)\qquad \varepsilon_2(\tau')=\varepsilon_3(\tau')=\deleted{\unsure{\kappa}}\chi(\tau')=-1.
\]
On a un merveilleux diagramme (p.\unsure{580'}).
\margin{664}


\scanpage{666}

\[
(58)\quad
\diagram{\begin{tikzcd}[ampersand replacement=\&,column sep=large,row sep=large]
\&D_\rho\arrow[r,no head]\&\underbrace{N_\rho\cap\unsure{\mathfrak{G}|\mathcal{G}}}_{=\mu_{\tau'}\cdot SZ_\rho}\arrow[rr,no head]\arrow[rr,bend left=18,no head,"a"]\arrow[dr,no head]\&\&N_\rho\\
D_{\tau'}=\mu_{\tau'}\times\mu.\arrow[ur]\arrow[dr,no head]\&\&\&\unsure{\mathfrak{G}|\mathcal{G}}_{\unsure{\rho}}\arrow[r,"\deleted{\unsure{\hookrightarrow}}" description]\&G_{\rho,\sigma}\arrow[u,bend left=12]\arrow[u,phantom,"\unsure{\nwarrow|\searrow}" description]\\
\&D_\sigma\arrow[r,no head]\&\underbrace{N_\sigma\cap\unsure{\mathfrak{G}|\mathcal{G}}}_{=\mu_{\tau'}\cdot SZ_\sigma}\arrow[rr,no head]\arrow[ur,no head]\&\&N_\sigma\arrow[u,no head]
\end{tikzcd}}
\]
ou mieux
\[
(59)\quad
\diagram{\begin{tikzcd}[ampersand replacement=\&,column sep=1.2em,row sep=1.6em]
\&\&\&D_\sigma\&\&\&\&\&G_{\rho,\sigma}\\
\mu_{\tau'}\&D_{\tau'}\&\&\&\&SN_\sigma\&\&N_\sigma\&\\
\&\&\&L_\sigma\&\&\&\unsure{\mathfrak{G}|\mathcal{G}}\&\&\\
\&\&\&\&\&SZ_\sigma\&\&Z_\sigma\&G_{\rho,\sigma}^{\prime\prime}\\
1\&\mu\&D_\rho\&\&SN_\rho\&\&\&N_\rho\&\\
\&\&\&\&\&\&\unsure{\mathfrak{G}|\mathcal{G}}^+\&\&\\
\&\&\&\&SZ_\rho\&\&\&\unsure{Z_\sigma|Z_\rho}\&G_{\rho,\sigma}'\\
\&\&L_\rho\&\&\&\&\&\&
\arrow[from=5-1,to=2-1,hook]
\arrow[from=5-1,to=5-2,hook]
\arrow[from=2-1,to=2-2,hook]
\arrow[from=5-2,to=2-2,hook]
\arrow[from=2-2,to=1-4,hook,"2"]
\arrow[from=5-2,to=3-4,hook,"2"]
\arrow[from=2-2,to=5-3,hook,"3"']
\arrow[from=5-2,to=8-3,hook,"3."']
\arrow[from=3-4,to=1-4,hook]
\arrow[from=8-3,to=5-3,hook]
\arrow[from=1-4,to=2-6]
\arrow[from=3-4,to=4-6,hook]
\arrow[from=5-3,to=5-5,hook]
\arrow[from=8-3,to=7-5,hook]
\arrow[from=4-6,to=2-6,hook]
\arrow[from=7-5,to=5-5,hook]
\arrow[from=2-6,to=3-7,hook]
\arrow[from=5-5,to=3-7,hook]
\arrow[from=4-6,to=6-7,hook]
\arrow[from=7-5,to=6-7,hook]
\arrow[from=6-7,to=3-7,hook]
\arrow[from=2-6,to=2-8,hook]
\arrow[from=5-5,to=5-8,hook]
\arrow[from=4-6,to=4-8,hook]
\arrow[from=7-5,to=7-8,hook]
\arrow[from=4-8,to=2-8,hook]
\arrow[from=7-8,to=5-8,hook]
\arrow[from=2-8,to=1-9,hook,bend right=12,"\unsure{b}"]
\arrow[from=3-7,to=1-9,hook,bend right=20]
\arrow[from=5-8,to=1-9,hook,bend right=15,"\deleted{\unsure{a_{\rho,\sigma}}}" description]
\arrow[from=4-8,to=4-9,hook]
\arrow[from=7-8,to=7-9,hook]
\arrow[from=4-9,to=1-9,hook]
\arrow[from=7-9,to=1-9,hook,bend right=10]
\arrow[from=6-7,to=4-9,dashed]
\arrow[from=6-7,to=7-9,dashed]
\end{tikzcd}}
\]

On définit :
\[
(60)\quad
N_\rho^{\unsure{\natural|\sharp}}\deleted{\unsure{=}\,\unsure{N_\rho\cap G_{\rho,\sigma}^{\natural}}}\quad
N_\sigma^{\unsure{\natural|\sharp}},Z_\rho^{\unsure{\natural|\sharp}},Z_\sigma^{\unsure{\natural|\sharp}},
\deleted{\unsure{SN_\rho^{\natural}},\unsure{SN_\sigma^{\natural}}},
SZ_\rho^{\unsure{\natural|\sharp}},SZ_\sigma^{\unsure{\natural|\sharp}}
\]
intersections des $N_\rho,N_\sigma$ etc avec $G_{\rho,\sigma}^{\unsure{\natural|\sharp}}$.

On a des suites exactes
\[
(61)\quad\left\{
\diagram{\begin{tikzcd}[ampersand replacement=\&]
1\arrow[r]\&S\unsure{Z|N}_\rho^{\unsure{\natural|\sharp}}\arrow[r]\&N_\rho^{\unsure{\natural|\sharp}}\arrow[r,"\delta_\rho^{\unsure{\natural|\sharp}}"]\&\gamma_{\rho,\sigma}\arrow[r,no head]\&1\\
1\arrow[r,no head]\&SZ_\sigma^{\unsure{\natural|\sharp}}\arrow[r]\&N_\sigma^{\unsure{\natural|\sharp}}\arrow[r,"\delta_\sigma^{\unsure{\natural|\sharp}}"]\&\gamma_{\rho,\sigma}\arrow[r,no head]\&1
\end{tikzcd}}\right.
\]
et des hom de suites exactes
\margin{665}


\scanpage{667}
\mspage{\unsure{536|636}}

\[
\tag{62}
\left\{\begin{gathered}
\diagram{\begin{tikzcd}[ampersand replacement=\&,column sep=large,row sep=large]
1\arrow[r] \& SZ_\rho^{\unsure{\natural|\sharp}}\arrow[r,no head]\arrow[d,hook] \& \unsure{N|Z}_\rho^{\unsure{\natural|\sharp}}\arrow[r,no head,"\delta_\rho^{\unsure{\natural|\sharp}}"]\arrow[d,hook] \& \gamma_{\rho,\sigma}\arrow[r,no head]\arrow[d,"\unsure{1|\Vert}"] \& 1\\
1\arrow[r,no head] \& S\unsure{Z|N}_\rho\arrow[r,no head] \& \unsure{N|Z}_\rho\arrow[r,"\delta_\rho"] \& \gamma_{\rho,\sigma}\arrow[r,no head] \& 1
\end{tikzcd}}\\[2ex]
\diagram{\begin{tikzcd}[ampersand replacement=\&,column sep=large,row sep=large]
1\arrow[r,no head] \& SZ_\sigma^{\unsure{\natural|\sharp}}\arrow[r,no head]\arrow[d,hook] \& N_\sigma^{\unsure{\natural|\sharp}}\arrow[r,no head,"\delta_\sigma^{\unsure{\natural|\sharp}}"]\arrow[d,hook] \& \gamma_{\rho,\sigma}\arrow[r,no head]\arrow[d,"\unsure{1|\Vert}"] \& 1\\
1\arrow[r,no head] \& SZ_\sigma\arrow[r,no head] \& N_\sigma\arrow[r,"\delta_\sigma"] \& \gamma_{\rho,\sigma}\arrow[r,no head] \& 1
\end{tikzcd}},
\end{gathered}\right.
\]

\[
\tag{63}
\left\{\begin{gathered}
\diagram{\begin{tikzcd}[ampersand replacement=\&,column sep=large,row sep=large]
1\arrow[r] \& \unsure{Z|I}_\rho^{\unsure{\natural|\sharp}}\arrow[r,no head]\arrow[d,"{\text{épi}}"] \& N_\rho^{\unsure{\natural|\sharp}}\arrow[r,"\varepsilon_{3\rho}^{\unsure{\natural|\sharp}}"]\arrow[d,"{\text{épi}}"] \& \mu\arrow[r,no head]\arrow[d,"\unsure{1|\Vert}"] \& 1\\
1\arrow[r] \& \gamma'_{\rho,\sigma}\arrow[r] \& \gamma_{\rho,\sigma}\arrow[r,"\varepsilon_{3\gamma}"] \& \mu\arrow[r,no head] \& 1
\end{tikzcd}},\\[2ex]
\diagram{\begin{tikzcd}[ampersand replacement=\&,column sep=large,row sep=large]
1\arrow[r] \& Z_\sigma^{\unsure{\natural|\sharp}}\arrow[r,no head]\arrow[d,"{\text{épi}}"] \& N_\sigma^{\unsure{\natural|\sharp}}\arrow[r,"\varepsilon_{2\sigma}^{\unsure{\natural|\sharp}}"]\arrow[d,"{\text{épi}}"] \& \mu\arrow[r,no head]\arrow[d,"\unsure{1|\Vert}"] \& 1\\
1\arrow[r,no head] \& \gamma''_{\rho,\sigma}\arrow[r] \& \gamma_{\rho,\sigma}\arrow[r,no head] \& \mu\arrow[r,no head] \& 1
\end{tikzcd}}
\end{gathered}\right.
\]

(et diagrammes analogues \deleted{\unsure{avec} \unsure{(63)}} \inserted{\unsure{sans}} $\unsure{\natural|\sharp}$).

\deleted{\unsure{d'où}} On a donc un diagramme

\[
\tag{64}
\diagram{\begin{tikzcd}[ampersand replacement=\&,column sep=small,row sep=small]
\& L_0^{\unsure{\natural|\sharp}}=SZ_{\rho,\sigma}\arrow[rrr,hook]\arrow[dddr,hook] \& \& \& Z_{\rho,\sigma}\arrow[dddr,hook] \& \& \& \mu \& \& \\
\& \& SZ_\rho^{\unsure{\natural|\sharp}}\arrow[rrr,hook]\arrow[dd,hook] \& \& \& N_\rho^{\unsure{\natural|\sharp}}\arrow[dd,hook] \& \& \& \widehat{\mathbb Z}^{*} \& \\
\& \& \& SZ_\sigma^{\unsure{\natural|\sharp}}\arrow[rrr,hook]\arrow[dl,hook] \& \& \& N_\sigma^{\unsure{\natural|\sharp}}\arrow[dl,hook] \& \& \& \mu\\
1\arrow[rr] \& \& \unsure{\mathfrak{G}|\mathcal{G}}^{\unsure{\natural|\sharp}}\arrow[rrr,no head]\arrow[ddl,hook]\arrow[ddrrr,no head,"{\deleted{\unsure{\hookrightarrow}}}" description] \& \& \& G_{\rho,\sigma}^{\unsure{\natural|\sharp}}\arrow[rr] \& \& \gamma_{\rho,\sigma}\arrow[rr,no head]\arrow[uuu,"\varepsilon_{3\gamma}"]\arrow[uur,"\chi_\gamma" description]\arrow[urr,"\varepsilon_{2\gamma}"'] \& \& 1\\
\& \& \& \& \& \& \& \& \& \\
1\arrow[r] \& \unsure{\mathfrak{G}|\mathcal{G}}^{+}\arrow[rrrr,hook] \& \& \& \& \unsure{\mathfrak{G}|\mathcal{G}}\arrow[rr]\arrow[uu,hook] \& \& \mu_\tau\arrow[rr]\arrow[uu,hook] \& \& 1
\end{tikzcd}}
\]

cf. p.581' (45)), et diagrammes analogues sans $\unsure{\natural|\sharp}$.

\margin{666}


\scanpage{668}

\deleted{\uline{Exemple \inserted{\unsure{du cas}} universel}
$\unsure{\mathfrak{G}|\mathcal{G}}=\unsure{\mathrm{GL}|\mathrm{SL}}(2,\mathbb Z)^\wedge$,
avec $\unsure{\mathfrak{G}|\mathcal{G}}^{\unsure{\natural|\sharp}}=\Pi_0$.}

Je voudrais maintenant définir l'analogue des \inserted{\unsure{ss-}}groupes
$\unsure{M|\mathcal{M}}(\tfrac12)$, $\unsure{M|\mathcal{M}}'(-\overline j)$,
$\unsure{N_\rho}^{?}$\inserted{$\unsure{\unsure{M|\mathcal{M}}'''(-\overline j)|\unsure{M|\mathcal{M}}''(-\overline j)}$}
de p.~615 ff,
ce qui implique qu'on ait un groupe $\subset G_{\rho,\sigma}^{\unsure{\natural|\sharp}}$
jouant le rôle de $\unsure{M|\mathcal{M}}'_0$ --- l'idéal serait que ce soit
$G_{\rho,\sigma}^{\unsure{\natural|\sharp}}$ lui-même.
Mais dans le cas universel
$\unsure{\mathfrak{G}|\mathcal{G}}^+=\unsure{\mathrm{SL}|\mathrm{GL}}(2,\mathbb Z)^\wedge$,
$\unsure{\mathfrak{G}|\mathcal{G}}^{\unsure{\natural|\sharp}}=\Pi_0$,
d'où
$Z_{\rho,\sigma}=\unsure{\unsure{M|\mathcal{M}}[0]'|\unsure{M|\mathcal{M}}[0]!}=\unsure{M|\mathcal{M}}(\unsure{0|\infty}).L_0^\pi$,
\inserted{$L_0^{\unsure{\natural|\sharp}}=L_0^\pi$}
(on trouve $G_{\rho,\sigma}\simeq \unsure{M|\mathcal{M}}$,
\deleted{et $G_{\rho,\sigma}={}$ss-\unsure{groupe}}
$G_{\rho,\sigma}^{\unsure{\natural|\sharp}}={}$ss-\unsure{groupe}
de $G_{\rho,\sigma}=\unsure{M|\mathcal{M}}$ engendré par
$\unsure{\mathfrak{G}|\mathcal{G}}^{\unsure{\natural|\sharp}}=\Pi_0$
et par $Z_{\rho,\sigma}=\unsure{M|\mathcal{M}}(\unsure{0|\infty})L_0^\pi$
$=\unsure{M|\mathcal{M}}(\unsure{0|\infty}).\Pi=\unsure{\unsure{M|\mathcal{M}}'|\unsure{M|\mathcal{M}}'!}$,
\uline{et non} $\unsure{M|\mathcal{M}}(\unsure{0|\infty}).\Pi_0=\unsure{M|\mathcal{M}}'_0$\deleted{$\unsure{!}$}

L'ennui provient finalement du fait que c'est $Z_{\rho,\sigma}$
qui est trop gros, on aimerait définir dans $Z_{\rho,\sigma}$ un sous-groupe
qui joue le rôle de $\unsure{M|\mathcal{M}}(\unsure{0|\infty}).\Pi_0$. Il n'y a pas de difficulté si on peut définir
dans \deleted{\unsure{le}} $Z_{\rho,\sigma}$ un sous-groupe qui
corresponde \unsure{à} $\unsure{M|\mathcal{M}}(\unsure{0|\infty})$
\[
Z_{\rho,\sigma}(0)\subset Z_{\rho,\sigma}
\tag{65}
\]
qui donne un scindage
\[
Z_{\rho,\sigma}\simeq Z_{\rho,\sigma}(\unsure{0|\infty}).L_0^{\unsure{\natural|\sharp}}
\qquad (\tfrac12\ \text{direct})
\tag{66}
\]
et on \deleted{\unsure{possèdera}}\inserted{\unsure{aura}} alors
\[
\deleted{\unsure{\unsure{M|\mathcal{M}}}_{\rho,\sigma}}\inserted{G_{\rho,\sigma}}
\simeq Z_{\rho,\sigma}(0).\unsure{\mathfrak{G}|\mathcal{G}}^+,
\qquad
\deleted{\unsure{G_{\rho,\sigma}^{\natural}}\simeq
Z_{\rho,\sigma}(0).\unsure{\mathfrak{G}|\mathcal{G}}^{\unsure{\natural|\sharp}}}
\qquad (\tfrac12\ \text{directs})
\tag{67}
\]
\margin{667}


\scanpage{669}
\mspage{637}

et on \deleted{\unsure{définira|définirait}} poserait
\[
\left\{
\begin{array}{l}
(G_{\rho,\sigma}^{\unsure{\natural|\sharp}})_0
\overset{\unsure{\mathrm{dfn}|\text{déf}}}{=}
Z_{\rho,\sigma}(\unsure{0|\infty}).
\unsure{\mathfrak{G}|\mathcal{G}}^{\unsure{\natural|\sharp}}
\subset G_{\rho,\sigma}^{\unsure{\natural|\sharp}}
\subset
\begin{array}{c}
G_{\rho,\sigma}\\[-2pt]
\Vert\\[-2pt]
Z_{\rho,\sigma}(\unsure{0|\infty}).\unsure{\mathfrak{G}|\mathcal{G}}^+
\end{array}
\\[10pt]
Z_{\rho,\sigma}^{\unsure{\natural|\sharp}}\unsure{\subset|=}
\deleted{\unsure{G}}Z_{\rho,\sigma}
\\[-2pt]
\qquad\Vert
\\[-2pt]
G_{\rho,\sigma\,\unsure{0}}^{\unsure{\natural|\sharp}}\cap Z_{\rho,\sigma}
=Z_{\rho,\sigma}(\unsure{0|\infty}).
(L_0\cap\unsure{\mathfrak{G}|\mathcal{G}}^{\unsure{\natural|\sharp}})
\end{array}
\right.
\tag{68}
\]
\margin{\deleted{\unsure{67}}.}

Rappelons une façon standard d'obtenir un scindage can. de $Z_{\rho,\sigma}$
sur $\gamma_{\rho,\sigma}$, quand on dispose d'un
\inserted{$(\rho,\sigma)$-}hom \deleted{$\unsure{\varphi}$}
\[
\diagram{\begin{tikzcd}[ampersand replacement=\&,column sep=large]
\unsure{\mathfrak{G}|\mathcal{G}}^+ \arrow[r]
\& \mathrm{SL}(2,\widehat{\mathbb Z})'
=\mathrm{SL}(2,\widehat{\mathbb Z})/\{\pm1\}
\end{tikzcd}}
\tag{69}
\]
commutant à l'action de $\tau$, ce qui induit un hom
\[
\diagram{\begin{tikzcd}[ampersand replacement=\&,column sep=large]
Z_{\rho,\sigma}\arrow[r]
\& B'_0=\left\{
\begin{pmatrix}\mu&\lambda\\0&1\end{pmatrix}
\;\middle|\; \mu\in\widehat{\mathbb Z}^*;\ \lambda\in\widehat{\mathbb Z}
\right\}
\end{tikzcd}}
\tag{70}
\]
(voir : l'étude du \S~48)

induisant \deleted{sur $L_0$}
\[
\left\{
\diagram{\begin{tikzcd}[ampersand replacement=\&,column sep=large,row sep=large]
L_0^{\unsure{\natural|\sharp}}\arrow[r]
\&L'_0=\left\{
\begin{pmatrix}1&\lambda\\0&1\end{pmatrix}
\;\middle|\;\lambda\in\widehat{\mathbb Z}
\right\}\\
\xi_0^n\arrow[r,mapsto]
\&\begin{pmatrix}\unsure{1}&n\\0&1\end{pmatrix}
\end{tikzcd}}
\right.
\tag{71}
\]
\deleted{\unsure{donc|d'où}}\inserted{plus précisément} on a un diagramme
de suites exactes
\[
\diagram{\begin{tikzcd}[ampersand replacement=\&,column sep=large,row sep=huge]
1\arrow[r]
\&L_0^{\unsure{\natural|\sharp}}\arrow[r]\arrow[d,"\xi_0^n\mapsto n"']
\&Z_{\rho,\sigma}\arrow[r]\arrow[d]
\&\gamma_{\rho,\sigma}\arrow[r,no head]\arrow[d,"\chi_{\unsure{\gamma|\sigma}}"]
\&1\\
1\arrow[r]
\&\widehat{\mathbb Z}\arrow[r,two heads,
  "{\unsure{\lambda\mapsto\left(\begin{smallmatrix}1&\lambda\\0&1\end{smallmatrix}\right)}}"']
\&B''_0\arrow[r,"\det",
  "{\left(\begin{smallmatrix}\mu&\lambda\\0&1\end{smallmatrix}\right)\mapsto\mu}"']
\&\widehat{\mathbb Z}^*\arrow[r,no head]
\&1
\end{tikzcd}}
\tag{72}
\]
Si l'hom (71) est un \uline{iso}, alors \unsure{tout}
\margin{668}


\scanpage{670}

scindage de la \inserted{\unsure{2ième|2ème}} suite exacte en définit par
image inverse un de la première, \deleted{\unsure{et le}}
Mais dans les cas \inserted{typiques} qui \unsure{nous},
$L_0^{\unsure{\natural|\sharp}}$ (qui n'est pas assez petit pour donner
"le bon" \deleted{$\unsure{\unsure{M|\mathcal{M}}}$} $G_{\rho,\sigma}^{\unsure{\natural|\sharp}}$ !)
n'est pas assez gros par contre pour donner un iso dans (71) !
Mais \deleted{\unsure{Rem|Alors}} on peut aussi faire une
\uline{hypothèse} supplémentaire --- savoir que l'image inverse par (70)
de \deleted{\unsure{la}}\inserted{"\unsure{bon|bonne}"} section
\[
T'_0=\left\{
\begin{pmatrix}\mu&0\\0&1\end{pmatrix}
\;\middle|\;\mu\in\widehat{\mathbb Z}^*
\right\}\subset B'_0
\quad\text{---}
\tag{73}
\]
qui est un sous-groupe
\[
\begin{array}{l}
\deleted{\unsure{\unsure{M|\mathcal{M}}}_{\rho,\sigma}}\inserted{Z_{\rho,\sigma}}
(\unsure{0|\infty})\subset Z_{\rho,\sigma}\\[-2pt]
\qquad\Vert\\[-2pt]
\left\{
u\in Z_{\rho,\sigma}\;\middle|\;
\begin{array}{l}
\text{son image }
\begin{pmatrix}\mu&\lambda\\0&1\end{pmatrix}
\text{ est dans }T'_0,\\
\hfill\text{i.e. }\lambda=0
\end{array}
\right\}
\end{array}
\tag{74}
\]
satisfaisant
\[
Z_{\rho,\sigma}(\unsure{0|\infty})\cap L_0^{\unsure{\mu|\natural}}
=\{1\},
\tag{75}
\]
soit en fait un sous-groupe section,
i.e. que $Z_{\rho,\sigma}(\unsure{0|\infty})\hookrightarrow\gamma_{\rho,\sigma}$
soit épi i.e. iso, i.e. que l'on a bien
\[
Z_{\rho,\sigma}
=Z_{\rho,\sigma}(\unsure{0|\infty}).L_0^{\unsure{\natural|\sharp}}.
\tag{76}
\]
J'ai des bonnes raisons de penser que (dès que (69) existe)
cette condition \unsure{est} satisfaite dans tous les cas vraiment utiles.
\deleted{quitte, s'il le faut, à}
\unsure{Ce|Le} choix d'un
\margin{669}


\scanpage{671}
\mspage{638}

$Z_{\rho,\sigma}(\unsure{0|\infty})$ a la bonne propriété d'ailleurs d'être
"fonctoriel" dans un sens évident (cf. plus bas les questions de fonctorialité),
ce qui sera essentiel pour avoir un bon formalisme. Ceci nous permet donc
de définir un sous groupe
$(G_{\rho,\sigma}^{\unsure{\natural|\sharp}})_0$
\inserted{$\subset G_{\rho,\sigma}^{\unsure{\natural|\sharp}}$}
par (68). Si sa définition n'était ainsi un peu trop alambi-\linebreak
quée, c'est
\deleted{\unsure{cette} définition de $\unsure{G}$}
\inserted{groupe \inserted{"plus fin"} qui serait}
\unsure{meilleur|meilleure},
et mériterait en fait la notation plus simple
$G_{\rho,\sigma}^{\unsure{\natural|\sharp}}$.

Autre façon de procéder !
\deleted{Dans le cas}\inserted{On cherche un}
\deleted{universel}
\unsure{système} \unsure{de|des} restrictions adéquates dans
la définition de $Z_{\rho,\sigma}$, qui
\deleted{\unsure{doivent}}
impliquent que
\inserted{$L_0^{\unsure{\natural|\sharp}}
=L_0\cap\unsure{\mathfrak{G}|\mathcal{G}}^{\unsure{\natural|\sharp}}$,
i.e. que}
pour $l\in L_0\unsure{\,a}$, on a
\[
(\mathrm{int}(l),1,1,1)\in Z_{\rho,\sigma}
\Longleftrightarrow
l\in\unsure{\mathfrak{G}|\mathcal{G}}^{\unsure{\natural|\sharp}}
\]
Mais dans le cas universel, avec
$\unsure{\mathfrak{G}|\mathcal{G}}^{\unsure{\natural|\sharp}}=\Pi_0$,
prenant $l=l_0$, son image dans
$\unsure{\mathfrak{G}|\mathcal{G}}^{\unsure{+|*}}/
\unsure{\mathfrak{G}|\mathcal{G}}^{\unsure{\natural|\sharp}}
=\unsure{\mathfrak{G}|\mathcal{G}}/\Pi_0$
est l'élément central non trivial de
$\unsure{\mathfrak{G}|\mathcal{G}}/\Pi_0$,
l'autom. int. qu'il définit dans
$\unsure{\mathfrak{G}|\mathcal{G}}^+/\Pi_0$
(et $\unsure{\overline{\mathrm m}|=}$ dans
$\unsure{\mathfrak{G}|\mathcal{G}}/\unsure{\pi|\tau}$)
est trivial --- donc ce n'est pas par
\deleted{\unsure{cette}} des propriétés de cette action
qu'on arrivera à l'\unsure{exorciser} !

Mais peut-être est-il prématuré d'essayer déjà de "découper" ces groupes
$\deleted{\unsure{\unsure{M|\mathcal{M}}}_{\rho,\sigma}^{\unsure{\natural|\sharp}}}
\inserted{G_{\rho,\sigma}^{\unsure{\natural|\sharp}}}$,
$\unsure{M|\mathcal{M}}_{\rho,\sigma}^{\unsure{\natural|\sharp}}$
\margin{670}


\scanpage{672}

"aussi fin que possible" --- peut-être y a-t-il
\inserted{encore} des erreurs d'orientation assez grossières qui
se rectifieront par l'examen des cas
particuliers. Donc je vais renoncer
pour l'instant à continuer à investir \deleted{\unsure{de}}
de l'énergie \unsure{au|en} découpage, et m'en
tiendrai donc là dans ce sens, quitte
à y revenir par la suite et à raffiner
les constructions présentes, qu'il faudra
cependant comparer à l'occasion à celles
du \S~48 $\overline{\underline{\unsure{\mathrm{XI}}}}$.

\margin{671}


\scanpage{673}
\mspage{639}

$5^\circ)\quad S_0\mathcal{M}_{\rho,\sigma}^{\unsure{\natural|\sharp}},\ \mathcal{M}_{\rho,\sigma}^{\unsure{\natural|\sharp}},\ S_0\Gamma_{\rho,\sigma}^{\unsure{\natural|\sharp}},\ \Gamma_{\rho,\sigma}\ \ldots .$

On posera comme dans p. 583
\[
\left\{\begin{aligned}
S_0\mathcal{M}_{\rho,\sigma}
 &= Z_{\rho,\sigma}\times_{\gamma_{\rho,\sigma}}N_\rho^{\unsure{\natural|\sharp}}
 \times_{\gamma_{\rho,\sigma}}N_\sigma^{\unsure{\natural|\sharp}}
 \times_{\gamma_{\rho,\sigma}}G_{\rho,\sigma}\\
S_0\mathcal{M}_{\rho,\sigma}^{\unsure{\natural|\sharp}}
 &= Z_{\rho,\sigma}\times_\gamma N_\rho^{\unsure{\natural|\sharp}}
 \times_\gamma N_\sigma^{\unsure{\natural|\sharp}}\times_\gamma G_{\rho,\sigma}^{\unsure{\natural|\sharp}}
\end{aligned}\right.
\tag{77}
\]

\diagram{\begin{tikzcd}[ampersand replacement=\&,column sep=small,row sep=large]
{\unsure{(78')|(78)}} \arrow[dd,bend right=70] \&
{\left\{\begin{aligned}
S_0\Gamma_{\rho,\sigma}&=Z_{\rho,\sigma}\times_\gamma N_\rho^{\unsure{\natural|\sharp}}\times N_\sigma^{\unsure{\natural|\sharp}}\\
\Gamma_{\rho,\sigma}&=\deleted{\unsure{Z_{\rho,\sigma}}}\ N_\rho^{\unsure{\natural|\sharp}}\times_\gamma N_\sigma^{\unsure{\natural|\sharp}}
\end{aligned}\right.}\\
(78) \&
{\left\{\begin{aligned}
\mathcal{M}_{\rho,\sigma}&\simeq N_\rho^{\unsure{\natural|\sharp}}\times N_\sigma^{\unsure{\natural|\sharp}}\times G_{\rho,\sigma}\\
\mathcal{M}_{\rho,\sigma}^{\unsure{\natural|\sharp}}&\simeq N_\rho^{\unsure{\natural|\sharp}}\times N_\sigma^{\unsure{\natural|\sharp}}\times G_{\rho,\sigma}^{\unsure{\natural|\sharp}}.
\end{aligned}\right.}\\
{} \& {}
\end{tikzcd}}

d'où diagramme de suites exactes

\[
\text{(79)}\qquad
\diagram{\begin{tikzcd}[ampersand replacement=\&,column sep=small,row sep=large]
\deleted{\unsure{1}}\ 1 \arrow[r] \&
\unsure{\mathfrak{G}|\mathcal{G}}^+ \arrow[rr,hook] \arrow[dr,hook] \&\&
S_0\mathcal{M}_{\rho,\sigma} \arrow[rr,hook] \arrow[dr,hook] \&\&
S_0\Gamma_{\rho,\sigma} \arrow[r] \arrow[dr,hook] \&1\&\\
\&1\arrow[r]\&\unsure{\mathfrak{G}|\mathcal{G}}^+\arrow[rr,hook]\&\&\mathcal{M}_{\rho,\sigma}\arrow[rr]\&\&\Gamma_{\rho,\sigma}\arrow[r]\&1\\
1\arrow[r]\&\unsure{\mathfrak{G}|\mathcal{G}}^{\unsure{\natural|\sharp}}\arrow[uu,hook]\arrow[rr,hook]\arrow[dr,"\unsure{\sim}"]\&\&
S_0\mathcal{M}_{\rho,\sigma}^{\unsure{\natural|\sharp}}\arrow[uu,hook]\arrow[rr]\arrow[dr,"{\deleted{\unsure{\nwarrow}}}"{pos=.1}]\&\&
S_0\Gamma_{\rho,\sigma}^{\unsure{\natural|\sharp}}\arrow[uu,hook]\arrow[dr,hook]\&\&\\
\&1\arrow[r]\&\unsure{\mathfrak{G}|\mathcal{G}}^{\unsure{\natural|\sharp}}\arrow[uu,hook]\arrow[rr,hook]\&\&
\deleted{\unsure{S_0}}\mathcal{M}_{\rho,\sigma}^{\unsure{\natural|\sharp}}\arrow[uu,hook]\arrow[rr]\&\&
\Gamma_{\rho,\sigma}^{\unsure{\natural|\sharp}}\arrow[uu,hook]\arrow[r]\&1
\end{tikzcd}}
\]

et des \inserted{\unsure{hom. de|hom de}} suites exactes
\[
\text{\unsure{(80)|(89)}}\quad
\diagram{\begin{tikzcd}[ampersand replacement=\&,column sep=small]
1\arrow[r]\&
{\overbrace{SZ_{\rho,\sigma}}^{L_0^{\unsure{\natural|\sharp}}}\times SN_\rho^{\unsure{\natural|\sharp}}\times SN_\sigma^{\unsure{\natural|\sharp}}\times \unsure{\mathfrak{G}|\mathcal{G}}^{\unsure{\natural|\sharp}}}
\arrow[r]\arrow[d]\&S_0\mathcal{M}_{\rho,\sigma}^{\unsure{\natural|\sharp}}\arrow[r]\arrow[d]\&\gamma_{\rho,\sigma}\arrow[r]\arrow[d,equal]\&1\\
1\arrow[r]\&
{\underbrace{SZ_{\rho,\sigma}}_{L_0^{\unsure{\natural|\sharp}}}\times SN_\rho^{\unsure{\natural|\sharp}}\times SN_\sigma^{\unsure{\natural|\sharp}}\times \unsure{\mathfrak{G}|\mathcal{G}}^+}
\arrow[r]\&S_0\mathcal{M}_{\rho,\sigma}\arrow[r]\&\gamma_{\rho,\sigma}\arrow[r]\&1
\end{tikzcd}}
\]
\[
\text{\unsure{(82)|(81)}}\quad
\diagram{\begin{tikzcd}[ampersand replacement=\&,column sep=small]
1\arrow[r]\&{SN_\rho^{\unsure{\natural|\sharp}}\times SN_\sigma^{\unsure{\natural|\sharp}}\times \unsure{\mathfrak{G}|\mathcal{G}}^{\unsure{\natural|\sharp}}}\arrow[r]\arrow[d]\&
\mathcal{M}_{\rho,\sigma}^{\unsure{\natural|\sharp}}\arrow[r]\arrow[d]\&\gamma_{\rho,\sigma}\arrow[r]\arrow[d,equal]\&1\\
1\arrow[r]\&{SN_\rho^{\unsure{\natural|\sharp}}\times SN_\sigma^{\unsure{\natural|\sharp}}\times \unsure{\mathfrak{G}|\mathcal{G}}}\arrow[r]\&
\mathcal{M}_{\rho,\sigma}\arrow[r]\&\gamma_{\rho,\sigma}\arrow[r]\&1
\end{tikzcd}}
\]
\[
\text{(83)}\quad
\diagram{\begin{tikzcd}[ampersand replacement=\&,column sep=small]
1\arrow[r]\&L_0^{\unsure{\natural|\sharp}}\times SN_\rho^{\unsure{\natural|\sharp}}\times SN_\sigma^{\unsure{\natural|\sharp}}\arrow[r]\&
S_0\Gamma_{\rho,\sigma}^{\unsure{\natural|\sharp}}\arrow[r]\&\gamma_{\rho,\sigma}\arrow[r]\&1
\end{tikzcd}}
\]
\[
\text{(84)}\quad
\diagram{\begin{tikzcd}[ampersand replacement=\&,column sep=small]
1\arrow[r]\&\deleted{\unsure{L_0^{\unsure{\natural|\sharp}}\times}}SN_\rho^{\unsure{\natural|\sharp}}\times SN_\sigma^{\unsure{\natural|\sharp}}\arrow[r]\&
\Gamma_{\rho,\sigma}\arrow[r]\&\gamma_{\rho,\sigma}\arrow[r]\&1
\end{tikzcd}}
\]
\margin{672}


\scanpage{674}

\[
\left\{\begin{aligned}
S_0\mathcal{M}_{\rho,\sigma}&\simeq\mathcal{M}_{\rho,\sigma}\times_{\gamma_{\rho,\sigma}} Z_{\rho,\sigma}\\
S_0\mathcal{M}_{\rho,\sigma}^{\unsure{\natural|\sharp}}&\simeq\mathcal{M}_{\rho,\sigma}^{\unsure{\natural|\sharp}}\times_{\gamma_{\rho,\sigma}}\unsure{Z_{\rho,\sigma}}
\end{aligned}\right.
\tag{85}
\]
\[
\deleted{\unsure{\{\,S_0}}\quad
S_0\Gamma_{\rho,\sigma}\simeq\Gamma_{\rho,\sigma}\times_{\gamma_{\rho,\sigma}}Z_{\rho,\sigma}
\tag{86}
\]
\[
\left\{\begin{aligned}
S_0\mathcal{M}_{\rho,\sigma}&\simeq\mathcal{M}_{\rho,\sigma}\times_{\Gamma_{\rho,\sigma}}S_0\Gamma_{\rho,\sigma}\\
S_0\mathcal{M}_{\rho,\sigma}^{\unsure{\natural|\sharp}}&\simeq\mathcal{M}_{\unsure{\rho|2},\sigma}^{\unsure{\natural|\sharp}}\times_{\Gamma_{\rho,\sigma}}S_0\Gamma_{\rho,\sigma}.
\end{aligned}\right.
\tag{87}
\]

\uline{\unsure{Rmq|Remq}} Cette \unsure{dégelée} de groupes \unsure{donne} un peu le mal de mer. Je ne suis pas sûr que les variantes $S_0$ $(S_0\mathcal{M},\ S_0\mathcal{M}^{\unsure{\natural|\sharp}},\ S_0\Gamma)$ soient très utiles.

Reste la \unsure{compétition|comparaison} entre des \unsure{invariants} avec ou sans $\unsure{\natural|\sharp}$, surtout \deleted{\unsure{les}} les $\mathcal{M}_{\rho,\sigma}$ et $\mathcal{M}_{\rho,\sigma}^{\unsure{\natural|\sharp}}$ --- dont la ``différence'' est ``la $\unsure{\underset{=}{\mathrm m}|\overline{\mathrm m}}$'' qu'entre $\unsure{\mathfrak{G}|\mathcal{G}}^+$ et $\unsure{\mathfrak{G}|\mathcal{G}}^{\unsure{\natural|\sharp}}$, ou \deleted{de façon \unsure{pré}} \inserted{et entre} $G_{\rho,\sigma}$ et $G_{\rho,\sigma}^{\unsure{\natural|\sharp}}$, de façon précise $\mathcal{M}_{\rho,\sigma}$ est invariant dans $\mathcal{M}_{\rho,\sigma}^{\unsure{\natural|\sharp}}$, et $G_{\rho,\sigma}^{\unsure{\natural|\sharp}}$ dans $G_{\rho,\sigma}$, on a des isomorphismes
\[
\text{(88)}\quad
\diagram{\begin{tikzcd}[ampersand replacement=\&,column sep=large]
\unsure{\mathfrak{G}|\mathcal{G}}^+/\unsure{\mathfrak{G}|\mathcal{G}}^{\unsure{\natural|\sharp}}\arrow[r,"\sim"]\&G_{\rho,\sigma}/G_{\rho,\sigma}^{\unsure{\natural|\sharp}}\arrow[r,"\sim"]\&\mathcal{M}_{\rho,\sigma}/\mathcal{M}_{\rho,\sigma}^{\unsure{\natural|\sharp}}.
\end{tikzcd}}
\]
Les groupes $G_{\rho,\sigma},\mathcal{M}_{\rho,\sigma}$ jouent \unsure{p. r. à|p. v. à} $G_{\rho,\sigma}^{\unsure{\natural|\sharp}},\mathcal{M}_{\rho,\sigma}^{\unsure{\natural|\sharp}}$ le rôle de fourre-touts (notamment pour nous permettre de considérer les autom. $\rho,\sigma$ de $G_{\rho,\sigma}^{\unsure{\natural|\sharp}},\mathcal{M}_{\rho,\sigma}^{\unsure{\natural|\sharp}}$ comme induits par des autom. int. des groupes
\margin{673}


\scanpage{675}
\mspage{640}

\unsure{environnants} $G_{\rho,\sigma},\mathcal{M}_{\rho,\sigma}$.

On a un élément
\[
\text{(89)}\quad
\begin{array}{cl}
\tau'_{\unsure{M|\mu}}\in\mathcal{M}_{\rho,\sigma}^{\unsure{\natural|\sharp}}
&\deleted{\text{\unsure{d'im.|d'où}}\ \tau'_\gamma}\ \inserted{\text{au dessus de }\tau'_\gamma\in\gamma_{\rho,\sigma}}\\
\Vert&\\[-1ex]
\diagram{\begin{tikzcd}[ampersand replacement=\&,column sep=small,row sep=large]
{\bigl(\unsure{\tau'_\rho},}\arrow[d,phantom,"{\in}"{rotate=90,description}]\&
{\tau'_\sigma,}\arrow[d,phantom,"{\in}"{rotate=90,description}]\&
{\tau'_G\bigr)}\arrow[d,phantom,"{\in}"{rotate=90,description}]\\
N_\rho^{\unsure{\natural|\sharp}}\&N_{\unsure{\rho|\sigma}}^{\unsure{\natural|\sharp}}\&
{\begin{array}{c}\deleted{\unsure{\mathfrak{G}}\in G_{\rho,\sigma}}\\\inserted{G_{\rho,\sigma}^{\unsure{\natural|\sharp}}}\end{array}}
\end{tikzcd}}&
\end{array}
\]

d'où \deleted{\unsure{des}} plongements canoniques
\[
\text{(90)}\quad
\diagram{\begin{tikzcd}[ampersand replacement=\&,column sep=large,row sep=large]
\unsure{\mathfrak{G}|\mathcal{G}}_\tau^{\unsure{\natural|\sharp}}\arrow[r,hook]\arrow[d,hook]\&\mathcal{M}_{\rho,\sigma}^{\unsure{\natural|\sharp}}\arrow[d]\&
{\begin{array}{l}\unsure{\mathfrak{G}|\mathcal{G}}^{\unsure{\natural|\sharp}}=\text{\unsure{image} inverse dans}\\
\qquad\mathcal{M}_{\rho,\sigma}\text{ de }\{1,\unsure{\tau_\gamma|\tau'_\gamma}\}\subset\gamma_{\rho,\sigma}\end{array}}\\
\unsure{\mathfrak{G}|\mathcal{G}}\arrow[r,hook]\&\mathcal{M}_{\rho,\sigma}\&
{\begin{array}{l}\unsure{\mathfrak{G}|\mathcal{G}}=\text{im. inv. dans }\mathcal{M}_{\rho,\sigma}\\
\qquad\text{de }\{1,\tau_\gamma\}\subset\gamma_{\rho,\sigma}\end{array}}
\end{tikzcd}}
\]

On peut reproduire les diagrammes \unsure{bordéliques} de p. \unsure{584}, avec des $\unsure{\natural|\sharp}$ à la clef\ldots

Je vais voir si je peux me dispenser justement de \unsure{redévelopper} les liens avec $\overline{Z}_{\rho,\sigma},\ S\unsure{G|g}_{\rho,\sigma}$ etc --- quitte à y revenir au besoin par la suite.
\margin{674}


\scanpage{676}

$6^\circ)$ \uline{Fonctorialités}.

\margin{\uline{NB} Jusqu'ici \inserted{(dans $1^\circ$ à $5^\circ$)} on n'avait pas utilisé l'hyp. que $\rho,\sigma$ \unsure{engendrent} $\unsure{\mathfrak{G}|\mathcal{G}}^+$}

Cf. p. 588. En plus de l'hom surjectif
\[
\text{(91)}\quad
\diagram{\begin{tikzcd}[ampersand replacement=\&]
\unsure{\mathfrak{G}|\mathcal{G}}\arrow[r,"\varphi"]\&\unsure{\mathfrak{G}|\mathcal{G}}'
\end{tikzcd}}
\]
\margin{Il vaut mieux \unsure{noter}
\diagram{\begin{tikzcd}[ampersand replacement=\&,column sep=small]\unsure{\mathfrak{G}|\mathcal{G}}_0\arrow[r]\&\unsure{\mathfrak{G}|\mathcal{G}}\end{tikzcd}}
cf p. $643'$ ff\ldots}

on suppose que $\varphi$ applique $\unsure{\mathfrak{G}|\mathcal{G}}^{\unsure{\natural|\sharp}}$ dans $\unsure{\mathfrak{G}|\mathcal{G}}^{\prime \unsure{\natural|\sharp}}$.

\deleted{Soit $u\in$} \inserted{On désigne ici par} $Z_{\rho,\sigma},\mathcal{M}_{\rho,\sigma}$ \unsure{et|etc} les invariants associés à $\unsure{\mathfrak{G}|\mathcal{G}}$, $Z_{\rho',\sigma'}$ etc ceux de $\unsure{\mathfrak{G}|\mathcal{G}}'$. On a des conditions équivalentes comme dans loc. cit --- où on va laisser tomber la considération de $\unsure{\Xi|\overline{Z}}_{\rho,\sigma}^{\unsure{\mathfrak{G}|\mathcal{G}}}$\ldots\ --- qui signifient que pour $\forall$
\[
\underline{u}=(u,\mu,\varepsilon_2,\varepsilon_3)\in Z_{\rho,\sigma}
\]
il existe un autom. $u'$ de $\unsure{\mathfrak{G}|\mathcal{G}}'$ rendant commutatif
\[
\text{(92)}\quad
\diagram{\begin{tikzcd}[ampersand replacement=\&,column sep=large,row sep=large]
\unsure{\mathfrak{G}|\mathcal{G}}^+\arrow[r,"\varphi"]\arrow[d,"u"']\&\unsure{\mathfrak{G}|\mathcal{G}}^{\prime+}\arrow[d,"u'"]\\
\unsure{\mathfrak{G}|\mathcal{G}}^{\unsure{++|+}}\arrow[r,"\varphi'"]\&\unsure{\mathfrak{G}|\mathcal{G}}^{\prime+}
\end{tikzcd}}
\]
lequel $u'$ sera unique. Il revient au $\unsure{=|\overline{\mathrm m}}$ de dire que les $u\in\operatorname{Aut}(\unsure{\mathfrak{G}|\mathcal{G}}^+)$ qui proviennent de $Z_{\rho,\sigma}$, i.e. qui \uline{normalisent} $L_0$, et qui transforment $\unsure{\xi_\rho|\ell_\rho},\unsure{\xi_\sigma|\ell_\sigma}$ en des \deleted{\unsure{éléments}}\inserted{\deleted{\unsure{groupes}}} \deleted{$\unsure{\mathfrak{G}|\mathcal{G}}^{\unsure{\natural|\sharp}}$-conjugués à} éléments $\unsure{\mathfrak{G}|\mathcal{G}}_\tau^{\unsure{\natural|\sharp}}$-conjugués à $\rho,\sigma$, stabilisent le noyau de $\varphi$ (91).

On trouve alors
\margin{675}


\scanpage{677}
\mspage{641}

\[
\text{(93)}\quad
\diagram{\begin{tikzcd}[ampersand replacement=\&]
Z_{\rho,\sigma}\arrow[r,"{\varphi_{\unsure{Z|z}}}"]\&Z_{\rho',\sigma'}
\end{tikzcd}}
\]
compatible avec les actions de $Z_{\rho,\sigma},Z_{\rho',\sigma'}$ sur $\unsure{\mathfrak{G}|\mathcal{G}}^+$ resp. $\unsure{\mathfrak{G}|\mathcal{G}}^{\prime+}$, et avec les hom
\[
\diagram{\begin{tikzcd}[ampersand replacement=\&,column sep=small,row sep=small]
L_0^{\unsure{\natural|\sharp}}\arrow[rr,bend right=35]\&{\subset Z_{\rho,\sigma}}\&{\unsure{\mathfrak{G}|\mathcal{G}}^+,}\&
L_0^{\prime \unsure{\natural|\sharp}}\arrow[rr,bend right=35]\&{\subset Z_{\rho',\sigma'}}\&{\unsure{\mathfrak{G}|\mathcal{G}}^{\prime+},}
\end{tikzcd}}
\]
d'où par la construction \unsure{Amal.} et un hom
\[
\text{(94)}\quad
\diagram{\begin{tikzcd}[ampersand replacement=\&]
G_{\rho,\sigma}\arrow[r,"\varphi_G"]\&G_{\rho',\sigma'}^{\unsure{\deleted{\unsure{\natural|\sharp}}|\varepsilon}}
\end{tikzcd}}
\]
induisant
\[
\text{(95)}\quad
\diagram{\begin{tikzcd}[ampersand replacement=\&]
G_{\rho,\sigma}^{\unsure{\natural|\sharp}}\arrow[r,"{\varphi_G^{\unsure{\natural|\sharp}}}"]\&G_{\rho',\sigma'}^{\unsure{\natural|\sharp}}
\end{tikzcd}}
\]
d'où un hom des diagrammes (64) relatif à \deleted{$\unsure{\mathfrak{G}|\mathcal{G}}$} $(\rho,\sigma)$, dans le diagramme analogue relatif à $(\rho',\sigma')$, avec ou sans $\unsure{\natural|\sharp}$.

Il n'en faut pas plus pour que toutes les constructions \inserted{et diagrammes} des n\textsuperscript{o} précédents, \unsure{relatifs|relatives} à $\rho,\sigma$, \unsure{s'envoient} dans celles relatives à $\rho',\sigma'$.

\uline{Remarque} Supposons que $\unsure{\mathfrak{G}|\mathcal{G}}'$ \inserted{$,\unsure{\mathfrak{G}|\mathcal{G}}^{\prime \unsure{\natural|\sharp}}$} \unsure{satisfasse} à la condition suivante: Pour $\unsure{\mathrm{tt}|\forall}$ couple
\[
\text{(96)}\quad
\left\{\begin{array}{l}
\text{d'éléments }\alpha',\beta'\in \unsure{\mathfrak{G}|\mathcal{G}}_\tau^{\prime \unsure{\natural|\sharp}},\quad\text{tels que l'on}\\[.6ex]
\text{ait}\\[.6ex]
\qquad[\alpha',\rho]=[\beta',\sigma]\,\unsure{\varepsilon_1|\xi_1}^{\mu-1}
\qquad\text{pour }\mu\in\widehat{\mathbb Z}\ \text{conv.}\\[.6ex]
\deleted{\text{(96)}}\\[-.5ex]
\text{il existe un }u'\in\operatorname{Aut}(\unsure{\mathfrak{G}|\mathcal{G}}^{\prime+})\quad\text{tel que}\\[.6ex]
\text{(97)}\qquad u'(\rho)=\alpha'(\rho),\quad u'(\sigma)=\beta'(\sigma)
\end{array}\right.
\]
i.e.
\[
\text{(98)}\qquad
\alpha^{\prime-1}u'\in\operatorname{Cent}_{\operatorname{Aut}}(\rho),\quad
\beta^{\prime-1}u'\in\operatorname{Cent}_{\operatorname{Aut}}(\sigma)
\]
\margin{676}


\scanpage{678}

la relation dans $\mathrm{Cent}$ (\unsure{96}) s'exprimant en termes de $u$ par la condition
\[
u'\bigl(L_0^{\unsure{\mathfrak{G}|\mathcal{G}}'}\bigr)
=L_0^{\unsure{\mathfrak{G}|\mathcal{G}}'}
\qquad\text{i.e. }u'\text{ normalise }\unsure{\mathfrak{G}|\mathcal{G}}'.
\tag{99}
\]
Alors \unsure{tout} \inserted{$(\rho,\sigma)$-}hom
$\bigl(\unsure{\mathfrak{G}|\mathcal{G}},\unsure{\mathfrak{G}|\mathcal{G}}^{\unsure{\natural|\sharp}}\bigr)
\unsure{\longrightarrow|-}
\bigl(\unsure{\mathfrak{G}|\mathcal{G}}',\unsure{\mathfrak{G}|\mathcal{G}}^{\prime\unsure{\natural|\sharp}}\bigr)$
est nécessairement admissible.

Nous allons appliquer ceci au cas \inserted{\unsure{où} $\unsure{\mathfrak{G}|\mathcal{G}}$ est} universel --- mais pour pouvoir en dire \deleted{\unsure{tout}} ce qu'on a envie, je vais d'abord parler des ss-groupes
\[
\mathcal{M}(\unsure{0|\infty}),\ \mathcal{M}(\rho),\ \mathcal{M}(\sigma)\subset\mathcal{M}_{\rho,\sigma}.
\]

\margin{677}


\scanpage{679}
\deleted{\unsure{646}}
\mspage{642}

\margin{$7^\circ)$ Les \unsure{ss-gr.}
$\mathcal{M}_{\rho,\sigma}^{\deleted{\unsure{\natural|\sharp}}}[\unsure{0|\rho}],\quad
\mathcal{M}_{\rho,\sigma}^{\deleted{\unsure{\natural|\sharp}}}(\rho),\quad
\mathcal{M}_{\rho,\sigma}^{\deleted{\unsure{\natural|\sharp}}}(\sigma).$}

Dans le groupe \deleted{$S_0\mathcal{M}_{\rho,\sigma}$}
\[
\begin{aligned}
S_0\mathcal{M}_{\rho,\sigma}^{\unsure{\natural|\sharp}}
&=Z_{\rho,\sigma}\times_{\gamma_{\rho,\sigma}}
N_\rho^{\unsure{\natural|\sharp}}\times_{\gamma_{\rho,\sigma}}
N_\sigma^{\unsure{\natural|\sharp}}\times_{\gamma_{\rho,\sigma}}
G_{\rho,\sigma}^{\unsure{\natural|\sharp}}\\
&=\left\{x=(u,a,b,U)\in Z_{\rho,\sigma}\times N_\rho\times N_\sigma\times G_{\rho,\sigma}
\;\middle|\;\delta_Z(u)=\delta_\rho(a)=\delta_\sigma(b)=\delta_G(U)\right\}
\end{aligned}
\tag{100}
\]
on introduit trois sous-groupes
\[
\left\{
\begin{array}{lll}
S_0\mathcal{M}_{\rho,\sigma}[0]\unsure{\subset|=}S_0\mathcal{M}_{\rho,\sigma}^{\unsure{\natural|\sharp}}
&\text{par condition}& U=i_Z(u)\\
S_0\mathcal{M}_{\rho,\sigma}[\rho]\subset S_0\mathcal{M}_{\rho,\sigma}^{\unsure{\natural|\sharp}}
&\text{\rule{2.6em}{0.3pt}}& U=i_\rho(a)\\
S_0\mathcal{M}_{\rho,\sigma}[\sigma]\subset S_0\mathcal{M}_{\rho,\sigma}^{\unsure{\natural|\sharp}}
&\text{\rule{2.6em}{0.3pt}}& U=i_\sigma(b)
\end{array}\right.
\tag{101}
\]
en utilisant maintenant les trois hom \unsure{cano.}
\diagram{
\[
\left\{\vcenter{\hbox{\begin{tikzcd}[ampersand replacement=\&,row sep=small,column sep=large]
i_Z:Z_{\rho,\sigma}\arrow[r,hook]\&G_{\rho,\sigma}\\
i_\rho:N_\rho\arrow[r,hook]\&G_{\rho,\sigma}\\
i_\sigma:N_\sigma\arrow[r,hook]\&G_{\rho,\sigma}
\end{tikzcd}}}\right.\;.
\tag{102}
\]
}

Ces trois groupes sont \unsure{manifestement} tous trois isomorphes :
$Z_{\rho,\sigma}\times_{\gamma_{\rho,\sigma}}
\deleted{\unsure{Z}}N_\rho^{\unsure{\natural|\sharp}}\times_{\gamma_{\rho,\sigma}}
N_\sigma^{\unsure{\natural|\sharp}}=S_0\Gamma_{\rho,\sigma}$,
de façon précise, dans la suite exacte (52)
\diagram{
\begin{tikzcd}[ampersand replacement=\&,column sep=large]
1\arrow[r]\&\unsure{\mathfrak{G}|\mathcal{G}}^+\arrow[r]\&S_0\mathcal{M}_{\rho,\sigma}\arrow[r]\&S_0\Gamma_{\rho,\sigma}\arrow[r]\&1
\end{tikzcd}
}
l'hom canonique
$S_0\mathcal{M}_{\rho,\sigma}\longrightarrow S_0\Gamma_{\rho,\sigma}$
induit des isom. entre les groupes (\unsure{68}) et $S_0\Gamma_{\rho,\sigma}$, les trois groupes sont donc des sous-groupes sections.

\margin{\deleted{(\unsure{69})}}
On définit $\mathcal{M}_{\rho,\sigma}[0],\ \mathcal{M}_{\rho,\sigma}(\rho)$
\inserted{(ou $\mathcal{M}_{\rho,\sigma}(\frac12)$)}
et $\mathcal{M}_{\rho,\sigma}(\sigma)$ (ou $\mathcal{M}_{\rho,\sigma}(-\bar{j})$) comme les images de
$S_0\mathcal{M}_{\rho,\sigma}[0],\ S_0\mathcal{M}_{\rho,\sigma}[\rho]$ et
$S_0\mathcal{M}_{\rho,\sigma}[\sigma]$ par l'hom canonique
$S_0\mathcal{M}_{\rho,\sigma}\longrightarrow\mathcal{M}_{\rho,\sigma}$, ce qui

\margin{678}


\scanpage{680}

revient à poser
\[
\left\{\begin{aligned}
\mathcal{M}_{\rho,\sigma}[0]
&=\{(a,b,U)\in\mathcal{M}_{\rho,\sigma}^{\unsure{\natural|\sharp}}\mid U\in Z_{\rho,\sigma}\}\\
\mathcal{M}_{\rho,\sigma}(\rho)
&=\{(a,b,U)\in\mathcal{M}_{\rho,\sigma}^{\unsure{\natural|\sharp}}\mid U=a\}\\
\mathcal{M}_{\rho,\sigma}(\sigma)
&=\{(a,b,U)\in\mathcal{M}_{\rho,\sigma}^{\unsure{\natural|\sharp}}\mid U=b\}.
\end{aligned}\right.
\tag{103}
\]

Donc les \unsure{sous-groupes|ss-groupes} $\mathcal{M}_{\rho,\sigma}(\rho),\mathcal{M}_{\rho,\sigma}(\sigma)$ de $\mathcal{M}_{\rho,\sigma}^{\unsure{\natural|\sharp}}$ sont des groupes sections au dessus de $\Gamma_{\rho,\sigma}^{\unsure{\natural|\sharp}}\simeq N_\rho^{\unsure{\natural|\sharp}}\mathbin{\times}_{\gamma_{\rho,\sigma}}N_\sigma^{\unsure{\natural|\sharp}}$, donc
\[
\left\{\begin{aligned}
\mathcal{M}_{\rho,\sigma}^{\unsure{\natural|\sharp}}
&\simeq\mathcal{M}_{\rho,\sigma}(\rho)\mathbin{.}\unsure{\mathfrak{G}|\mathcal{G}}^{\unsure{\natural|\sharp}}\\
\mathcal{M}_{\rho,\sigma}^{\unsure{\natural|\sharp}}
&\simeq\mathcal{M}_{\rho,\sigma}(\sigma)\mathbin{.}\unsure{\mathfrak{G}|\mathcal{G}}^{\unsure{\natural|\sharp}}
\end{aligned}\right.
\qquad(\tfrac12\ \text{\unsure{directs|direct=}}),
\tag{104}
\]

et de $\unsure{\overline{\mathrm m}|=}$
\[
\left\{\begin{aligned}
\mathcal{M}_{\rho,\sigma}&\simeq\mathcal{M}_{\rho,\sigma}(\rho)\mathbin{.}\unsure{\mathfrak{G}|\mathcal{G}}^+\\
\mathcal{M}_{\rho,\sigma}&\simeq\mathcal{M}_{\rho,\sigma}(\sigma)\mathbin{.}\unsure{\mathfrak{G}|\mathcal{G}}^+
\end{aligned}\right.
\qquad(\tfrac12\ \text{directs}).
\tag{105}
\]

D'autre part, on \unsure{aura}
\diagram{
\[
\begin{tikzcd}[ampersand replacement=\&,column sep=small,row sep=large]
1\&\unsure{\longrightarrow|-}\&L_0^{\unsure{\natural|\sharp}}\arrow[dd,"1"]\&-\&
\mathcal{M}_{\rho,\sigma}^{\deleted{\unsure{\natural|\sharp}}}[0]\&\longrightarrow\&\Gamma_{\rho,\sigma}\&-\&1\\
\&\&\&\&S_0\mathcal{M}_{\rho,\sigma}[0]\arrow[u,"\unsure{\simeq|\sim}"]\arrow[d,"\unsure{\simeq|\sim}"']\&\&\&\&\\
1\&-\&L_0^{\unsure{\natural|\sharp}}\&\longrightarrow\&S_0\Gamma_{\rho,\sigma}\&\unsure{\longrightarrow|-}\&\Gamma_{\rho,\sigma}\arrow[uu,"1"']\&-\&1
\end{tikzcd},
\tag{106}
\]
}

donc l'extension $\mathcal{M}_{\rho,\sigma}[0]$ de $\Gamma_{\rho,\sigma}$ par $L_0^{\unsure{\natural|\sharp}}$ \unsure{s'identifie} à l'ext. $S_0\Gamma_{\rho,\sigma}$ de $\Gamma_{\rho,\sigma}$ par $L_0^{\unsure{\natural|\sharp}}$, i.e. à l'image inverse de l'extension $Z_{\rho,\sigma}$ de $\gamma_{\rho,\sigma}$ par $L_0^{\unsure{\natural|\sharp}}$, via l'hom canonique

\margin{679}


\scanpage{681}
\mspage{643}

\[
\Gamma_{\rho,\sigma}
=N_\rho^{\unsure{\natural|\sharp}}\mathbin{\times}_{\gamma_{\rho,\sigma}}N_\sigma^{\unsure{\natural|\sharp}}
\longrightarrow\gamma_{\rho,\sigma}.
\tag{107}
\]

Donc la donnée d'une section de $\deleted{\unsure{\Gamma}}Z_{\rho,\sigma}$ sur $\gamma_{\rho,\sigma}$ (cf à ce sujet p. $636'$ ff)\unsure{.} en définit une de $S_0\Gamma_{\rho,\sigma}$ sur $\Gamma_{\rho,\sigma}$, donc de $\mathcal{M}_{\rho,\sigma}[0]$ sur $\Gamma_{\rho,\sigma}$. Quand une telle section est fixée, on \unsure{dénote} \unsure{ce} groupe section par
\[
\mathcal{M}_{\rho,\sigma}(\unsure{0|\infty})
\subset\mathcal{M}_{\rho,\sigma}^{\unsure{\natural|\sharp}},
\]

donc on aura
\[
\left\{\begin{aligned}
\mathcal{M}_{\rho,\sigma}[0]
&\simeq\mathcal{M}_{\rho,\sigma}(\unsure{0|\infty})\mathbin{.}L_0^{\unsure{\natural|\sharp}}\\
\mathcal{M}_{\rho,\sigma}^{\unsure{\natural|\sharp}}
&\unsure{\mathbin{.}}\simeq\mathcal{M}_{\rho,\sigma}(\unsure{0|\infty})\mathbin{.}\unsure{\mathfrak{G}|\mathcal{G}}^{\unsure{\natural|\sharp}}\\
\mathcal{M}_{\rho,\sigma}
&\simeq\mathcal{M}_{\rho,\sigma}(0)\mathbin{.}\unsure{\mathfrak{G}|\mathcal{G}}
\end{aligned}\right.
\;.
\tag{108}
\]

\margin{680}
